% =============================================================
%  chapterB_body.tex -- Appendix B: literature review of the Bushveld Complex
%
%  Moved verbatim on 2026-10-27 from Chapter 5's "Geological Background"
%  (Chapter5/Writing/chapter5_body.tex; pre-move copy
%  chapter5_body.tex.pre_litreview_move_20261027) by
%  Chapter5/Workspace/move_litreview_to_appendix_20261027.py. No sentence was
%  changed. Editorial changes only:
%    - heading levels promoted one step (the review is now its own chapter);
%    - sections reordered: Layered Mafic Intrusions -> The Bushveld Event and
%      RLS -> Structure and Relationship to Country Rocks -> Isotopic and
%      Volatile Constraints (NEW grouping heading over the radiogenic-isotope,
%      stable-isotope, and volatile-phase subsections, previously nested under
%      "The Bushveld Event and RLS") -> Petrological Debates, which now closes
%      with The Transformation Hypothesis (previously under Structure);
%    - \label{sec:ch5-...} names kept so existing cross-references still work.
% =============================================================

The Paleoproterozoic Rustenburg Layered Series (RLS) of the Bushveld
Magmatic Province in northeast South Africa is the largest layered
igneous intrusion in the world, and contains the world's largest source
of platinum group elements as well as important reserves of chromium,
vanadium, titanium, and iron. It is exceptionally well-preserved, and
has been the subject of intense study and debate for generations.
Studies focusing on the intrusion's margin have been especially fruitful
in informing geologists' understanding of igneous petrogenesis and
intrusion mechanisms, modes, and tempos of layered mafic intrusions and
plutonic rocks more generally. 

\section{Layered Mafic Intrusions}
\label{sec:ch5-layered-mafic-intrusions}

Layered mafic intrusions (LMIs) are compositionally zoned formations of
mafic rocks characterized by distinct millimeter- to meter-scale layering
of igneous minerals. Traditionally, the layering and chemical zonation
of these geologic marvels have been attributed to processes such as
fractional crystallization and gravitational settling within a magma
chamber \citep{wager1960,charlier2015}. Notable examples of
LMIs include the Stillwater (2.71 Ga) in Montana, USA; Great Dyke (2.57
Ga) in central Zimbabwe; the Bushveld's layered series (2.06 Ga) in
South Africa\textquotesingle s Limpopo, Mpumalanga, and Gauteng
provinces; the Muskox intrusion (1.27 Ga) in Nunavut, Canada; the Duluth
Complex (1.01 Ga) in Minnesota, USA; the Sept Iles intrusion (565 Ma) in
Quebec, Canada; Dufek (189 Ma) in Antarctica; the Rum intrusion (60.5
Ma) in the Inner Hebrides, Scotland; and the Skaergaard intrusion (55.4
Ma) in East Greenland. Despite their significance as supposedly simple
natural laboratories for studying the chemical evolution of magmas,
numerous aspects of their formation remain subject to debate, even
including the presence of significant liquid magma in their formation
\citep{kruger2020,mungall2016,latypovetal2022,yaoetal2021}.

Experimental observations have revealed that plagioclase flotation
occurs in mafic magmas at depths corresponding to the
Earth\textquotesingle s crust \citep{campbell1978}. This challenges
the notion that gravitational settling alone is responsible for igneous
layering. Instead, chemical means of phase separation are often
suggested, referencing processes like periodic nucleation during in-situ
crystallization \citep{mcbirney1979}, in-situ crystallization of a
magma with compositional layering developed and sustained through
double-diffusive convection \citep{huppert1981}, or preferential
dissolution and recrystallization of existing unstructured crystal
cumulates or mushes to minimize crystal surface energy perpendicular to
temperature gradients \citep{boudreau2011} or by preferential attack by
chlorine-rich fluids \citep{sonnenthal1992}. An alternative mechanical
explanation, seismically aided hydrodynamic sorting within a
tilting, crystal-rich cumulate pile, has also enjoyed support \citep{maier2013}.

\section{The Bushveld Event and RLS}
\label{sec:ch5-the-bushveld-event-and-rls}

The Rustenburg Layered Series (RLS) is renowned as the
world's largest layered mafic intrusion, distinguished
by its remarkably continuous layers, sometimes across hundreds of
kilometers \citep{cawthorn2015,latypovetal2022,scoon1994}. Within the RLS, pegmatoid ``reefs'' and chromitites host the
world's largest reserves of platinum group elements or
PGEs \citep{lee1996}. In addition to the RLS, the Bushveld event
encompasses several satellite intrusions, including the Uitkomst
Intrusion, Phalaborwa Carbonatite, the Stavoren Granophyre at the top of
the RLS, the massive felsic Rooiberg lava eruption covering 200,000~km\textsuperscript{3},
and the world's largest A-type granitic pluton at the
intrusion's roof: the Lebowa granite suite \citep{hill1996,kleemann1989,scoates2021,walraven1985}.

The mafic plutonic rocks of the RLS are subdivided into the Lower Zone,
Critical Zone, Main Zone, and Upper Zone, with unlayered sills and chill
sequences at the margin known as the Marginal Zone \citep{vongruenewaldt1985}. These zones are primarily distinguished by their mineralogy,
which transitions from more primitive assemblages in the Lower Zone to
more evolved ones in the Upper Zone. The Lower Zone uniquely features
dunites and harzburgites, the Critical Zone features chromitites and
norites, the Main Zone features gabbronorites, and lastly the Upper Zone
features diorites and magnetitites. The RLS exhibits significant
compositional zoning, with increasing iron, silica, vanadium, and alkali
content as well as chromium and magnesium depletion toward the
intrusion\textquotesingle s roof \citep{cawthorn2015}. While the RLS often
maintains remarkable lateral continuity on the kilometer scale,
variations in stratigraphy, PGE mineralization, and thickness sometimes
present challenges to stratigraphic correlations \citep{ashwal2005,mcdonald2005,scoon1994}.

\section{Structure and Relationship to Country Rocks}
\label{sec:ch5-structure-and-relationship-to-country-rocks}

The layered series outcrops as distinct lobes on the present-day
surface, with variations in stratigraphy, PGE mineralization, and
mineralogy observed across limbs. These lobes are separated and
influenced by geological features like the Thabazimbi-Murchison and
Steelpoort lineaments \citep{duplessis1990,ashwal2005,mcdonald2005,scoon1994}. Despite these
variations, there is remarkable continuity between layers. Chromitite
layers with similar thickness and chemical signatures can be correlated
across limbs separated by hundreds of kilometers of felsic igneous rocks
and sedimentary cover. Additionally, crustal thickness models derived
from seismic receiver functions and gravity anomalies support the
possibility that the intrusion is connected at depth \citep{webb2004,cawthorn2001}. Some researchers have further classified
the intrusion into ``distal'' and ``proximal'' facies based on the
hypothesis that the intrusion\textquotesingle s feeders exploited the
Thabazimbi-Murchison Lineament, supported by observations such as the
best development of ultramafic zones closest to this lineament \citep{scoon1994,clarke2000}.

The RLS generally dips inward toward its center, with the most
significant dips occurring at the margin \citep{button1976,frick1967}.
While the RLS was historically considered a lopolith \citep{wilson1956},
paleomagnetic studies favor an initially horizontal orientation that
flexed isostatically after crystallization ceased \citep{hattingh1986}. The
floor of the RLS exhibits synmagmatic deformation, possibly resulting
from interactions during the coalescence of offset interfingered sills
\citep{pollard1975,clarke2009}. Abundant large sedimentary
inliers within the complex can also be explained by the inhomogeneous
propagation of finger sills along the complex\textquotesingle s margin
\citep{clarke2000,clarke2009}.

\section{Isotopic and Volatile Constraints on the RLS}
\label{sec:appB-isotopic-volatile-constraints}

\subsection{Radiogenic Isotope Stratigraphy of the RLS}
\label{sec:ch5-radiogenic-isotope-stratigraphy-of-the-rls}

\citet{kruger1994} reported a comprehensive strontium isotope stratigraphy of
the RLS, revealing radiogenic Sr isotopes indicative of varying degrees
of crustal contamination throughout the intrusion. This stratigraphy
differentiates the Lower, Critical, and lower Main Zones from the upper
Main Zone and Upper Zone. The lower portion of the intrusion records
highly variable Sr isotopes, interpreted to arise due to repeated magma
replenishments in an open system. \citet{kruger1994} refers to this zone as
the ``Integration Stage.'' Conversely, the middle Main Zone displays
fairly consistent isotopic compositions, transitioning to more primitive
values at the Pyroxenite Marker---a well-known bed in the upper Main
Zone signifying a shift towards more primitive magmas \citep{cawthorn1991}. These two stratigraphic regions with internally constant initial
strontium isotope ratios are referred to as the ``Differentiation
Stage'' and could possibly have evolved through closed-system fractional
crystallization on the basis of strontium isotopes. Studies on neodymium
isotopes within the RLS are less detailed than those available for
strontium isotopes, but show a broadly similar pattern that requires
crustal contamination throughout the entire complex, with the Main and
Upper Zones requiring more radiogenic (crustal?) contamination than the
Lower and Critical Zones \citep{maier2000}.

Strontium isotope disequilibrium is well-known within the Merensky and
Bastard reefs, where pyroxenes and plagioclase have different and
variable initial strontium isotope values \citep{seabrook2005,cawthorn2015};
however, strontium isotope disequilibrium is not unique
to the Merensky and Bastard reefs. \citet{chutas2012} demonstrated
that plagioclase throughout the Lower and Critical Zones was isotopically
zoned with primitive cores and evolved rims, and that coexisting
pyroxenes and plagioclase were not in isotopic equilibrium. \citet{roelofse2012}
similarly observed significant strontium isotopic
disequilibrium in lower Main Zone rocks. This means that crystalline
phases from the entire Lower, Critical, and lower Main Zone cumulates
interacted with exotic magmas or hydrothermal fluids after substantial
crystallization, either by later intrusion of an exotic magma into a
preexisting crystal mush at the level of the RLS \citep{chutas2012}
or by intrusion of xenocryst-bearing crystal-rich magmas derived from a
subcompartmentalized chamber at depth \citep{roelofse2012}. Any
rock in the Critical, Lower, and lower Main Zones was thus formed by
multiple magmatic ± hydrothermal events/environments, and any bulk
isotopic or chemical measurement of these zones will inherently
integrate over multiple crystallization episodes.

A comprehensive dataset of rhenium-osmium (Re-Os) isotopes provides
valuable insights into the history of platinum group elements (PGEs)
within the Bushveld Complex. \citet{schoenberg1999} measured a 2.043 ±
0.011 Ga isochron age for the Bastard Reef, which contained little
common osmium. Furthermore, they report a systematic trend in the
initial \textsuperscript{187}Os/\textsuperscript{188}Os ratios within
chromitites spanning the Critical Zone. These ratios exhibit a distinct
progression, starting from values closely resembling chondritic
compositions and escalating towards highly radiogenic values as one
moves upwards through the stratigraphic sequence reaching their most
radiogenic values at the Merensky Reef and upper group (UG) chromitites
\citep{schoenberg1999,mccandless1999}.

The presence of heterogeneous initial
\textsuperscript{187}Os/\textsuperscript{188}Os values within
chromitites is commonly invoked to support the argument for an
orthomagmatic origin for PGE mineralization in the upper Critical Zone,
arguing against a hydrothermal origin \citep{schoenberg1999,mccandless1999}. However, the validity of the argument depends
upon whether Re-Os isotopes faithfully record magmatic processes.
Notably, experimental evidence has demonstrated that osmium diffuses
rapidly within sulfides, displaying a closure temperature similar to the
closure temperature of argon within biotite \citep{brenan2000}. Ar-Ar
thermochronology of biotite therefore can be used in this case as a
proxy to determine the timing of the closure of the Re-Os system. \citet{nomade2004}
reported an average Ar-Ar age of 2.042 ± 0.003 Ga for
biotite within the UG2 chromite, which agrees well with the Re-Os
isochron age of the Bastard Reef. Furthermore, this congruence extends
beyond the Bushveld Complex, as similar Ar-Ar ages are observed on a
regional scale within syn-kinematic micas found within low-grade
phyllitic metasediments of the Transvaal Supergroup \citep{alexandre2006}. These Ar-Ar ages effectively date regional compression within the
Kaapvaal craton, including deformation into fold and thrust belts
\citep{alexandre2006}. Consequently, it can be inferred that the
Re-Os isotope system underwent diffusive reequilibration within the
chromitites approximately 10 million years after the cessation of
magmatism within the RLS, which ended at 2.055 Ga \citep[age derived from][]{scoates2021,zeh2015}.

The most enigmatic aspect of the Re-Os isotope geochemistry of the RLS
centers around the coexistence of two unusual, isotopically distinct
sulfide populations within the Merensky Reef. A substantial population
of laurite (RuS\textsubscript{2}) exhibits the most radiogenic initial
\textsuperscript{187}Os/\textsuperscript{188}Os values of the entire 8-kilometer
stratigraphy of the RLS, while the same hand sample contains
erlichmanite (OsS\textsubscript{2}) with an osmium isotopic ratio
mirroring a chondritic value, the least radiogenic
\textsuperscript{187}Os/\textsuperscript{188}Os within the entire RLS
\citep{hart1989}. It is difficult to imagine how sulfides with
such different \textsuperscript{187}Os/\textsuperscript{188}Os values
were able to coexist within the same rock without isotopically
equilibrating.

In a broader geological context, \citet{richardson2008} conducted
measurements of sulfide Re-Os isotopes, as well as silicate Rb-Sr and
Sm-Nd isotopes, within kimberlitic diamond inclusions in close proximity
to the Bushveld Complex. Their aim was to constrain the isotopic
reservoirs within the subcontinental lithospheric mantle (SCLM) beneath
the Bushveld Complex. Their findings reveal the enrichment of radiogenic
isotopes across all three isotopic systems within the SCLM, which they
use to argue that the radiogenic isotopic signatures observed within the
RLS can be explained as a three-part mixture of mantle reservoirs that
requires no contributions from the crust, instead requiring only the
harzburgitic SCLM, eclogitic SCLM, and the depleted, convecting mantle.

\subsection{Stable Isotopes of the RLS}
\label{sec:ch5-stable-isotopes-of-the-rls}

Oxygen isotopic studies show that oxygen within the RLS was equilibrated
at magmatic temperatures, remarkably unaffected by interactions with
surface-derived fluids \citep{schiffries1989,harris2005}.
Early studies indicated magmatic temperatures for the Upper Zone
equilibrating at 1000°C, with the rest of the intrusion at 1200°C
\citep{schiffries1989}. However, later studies produced scattered
equilibration temperatures as low as 550°C \citep{harris2005}. Oxygen
isotopes within the RLS exhibit remarkable consistency throughout the
intrusion and are slightly heavier (\textasciitilde8‰
δ\textsuperscript{18}O) than mantle values (\textasciitilde6‰
δ\textsuperscript{18}O) suggesting substantial but well-mixed
contamination from surface-derived material, possibly from Transvaal
Supergroup rocks exposed on the surface or from unexposed, isotopically
heavy rocks at depth \citep{schiffries1989,harris2005}.
Similarly, multiple-sulfur isotopes identified a well-mixed and uniform
non-zero mass independent fractionation (MIF) of Δ\textsuperscript{33}S
\citep{pennistondorland2012} across the RLS, requiring a component
of surface-derived sulfur that was very efficiently mixed throughout the
entirety of the RLS. The remarkable homogeneity in volatile element
isotopes throughout the RLS contrasts with heterogeneous strontium and
neodymium isotopes throughout the intrusion, alluding to a decoupling of
magmatic and hydrothermal material sources within the RLS. The volatile
element isotopic data require a single source for late-stage
hydrothermal fluids across the complex.

Hydrogen isotopes indicate late-stage hydrothermal metasomatism at low
water/rock ratios \citep{schiffries1990}. Carbonated fluid within
the RLS imparted a carbon isotope signature likely arising from the
devolatilization of Transvaal carbonates \citep{schiffries1990}.
Sedimentary country rocks display oxygen and carbon isotopic depletions
correlating with increasing metamorphic grade and proximity to the RLS
margins \citep{buick2000,schiffries1990}. The lack of
complete isotopic equilibration of metasedimentary country rocks and RLS
fluids suggests that fluids flowed from devolatilized metasedimentary
rocks into the RLS, and not vice versa \citep{buick2000}.

\subsection{Unusual Volatile-bearing Phases in the Lower RLS}
\label{sec:ch5-unusual-volatile-phases-in-the-lower-rls}

Within the Lower and Critical Zones of the RLS, volatile-bearing phases
such as apatite, amphibole, and phlogopite are rare but often contain
exceptionally high chlorine concentrations \citep{boudreau1986,willmore2000}. These phases are most concentrated in the
lateral Merensky Reef and the vertical metasomatic Driekop dunite pipe,
often associated with magmatic graphite and platiniferous sulfides
\citep{ballhaus1985,schiffries1982}. Thermodynamic
calculations suggest that the initial volatile phase exsolved from RLS
magmas was predominantly hydrochloric acid (HCl) and chloromethane
(CH\textsubscript{3}Cl), with minor amounts of water and sulfur \citep{mathez1989}.
This fluid likely exsolved sulfur during graphite crystallization.
Chlorine isotopic studies of the RLS show that these Cl-rich fluids were
not derived from the local Transvaal Supergroup, but are consistent with
an origin from SCLM charged with volatiles from dehydrating subducting
altered oceanic crust \citep{willmore2002}. The association of these
C-Cl rich volatile phases with platiniferous horizons has led many
authors to suggest that PGE mineralization is controlled by C-Cl rich
fluids, perhaps by chromatographic extraction by reactive porous flow of
C-Cl rich fluids within the existing Critical Zone cumulates \citep{boudreau1999}.

\section{Petrological Debates Within the RLS}
\label{sec:ch5-petrological-debates-within-the-rls}

\subsection{Open or Closed System}
\label{sec:ch5-open-or-closed-system}

Another quintessential LMI, the Skaergaard intrusion, is thought to have
formed primarily through closed-system fractional crystallization from
the chamber\textquotesingle s outer regions inward \citep{mcbirney1979},
albeit with strontium isotopic evidence of country rock
assimilation in its early stages \citep{stewart1990}. In
contrast, the RLS is interpreted as having crystallized mainly in an
open system, featuring frequent magmatic recharge. This is supported by
numerous chemical reversals in plagioclase and pyroxene compositions, as
well as variations in initial strontium isotopes within the Lower,
Critical, and lower Main Zones \citep{kruger1994,cawthorn2015}. However,
the upper Main and Upper Zones (sometimes lumped together as the UUMZ)
of the RLS are often seen as products of closed-system fractional
crystallization, as indicated by uniform strontium isotopes \citep{kruger1994},
oxygen isotopes that seem to have evolved through closed-system
fractional crystallization \citep{bucholz2017}, and well-fitting
crystallization sequences produced by MELTS simulations of closed-system
fractional crystallization of plausible UUMZ parental compositions
\citep{vantongeren2010}. Nevertheless, this view has been challenged
by the identification of significant reversals in plagioclase and
nonconformable field relationships within the UUMZ \citep{yuan2017,scoon2012}.

\subsection{Spatial and Temporal Scales of Magma Chambers in the Crust, Pluton Assembly Mechanisms}
\label{sec:ch5-magma-chamber-scales}

The RLS, along with other LMIs, has become a battleground for debates
regarding the temporal and spatial scales of crustal magma chambers and
pluton assembly. Some argue that plutonic rocks represent fossilized
magma chambers \citep{kruger2020,bachmann2004,paterson2009}, while others propose that these rocks are assembled
incrementally through sill-like intrusions. Such arguments are applied
to rocks within the RLS by \citet{mitchell2019}, \citet{yaoetal2021},
\citet{scoon2023}, and \citet{bourdeau2023}. The latter
viewpoint first gained significant support when detailed, high-precision
zircon U/Pb geochronology studies revealed wide ranges of
crystallization ages within seemingly homogeneous, continuous plutons
\citep{coleman2004,glazner2004}. These wide age ranges
produced by zircon U/Pb geochronology in plutons greatly exceed the
timeframes over which large magma chambers can be thermally sustained
before cooling \citep{annen2009,cawthornwalraven1998}.

Zircons from the RLS and Stillwater layered mafic intrusions similarly
yield some out-of-sequence U/Pb ages that challenge the idea that they
represent large, fossilized magma chambers \citep{mungall2016,wall2018,scoates2021}. However, resolvably out-of-sequence
zircons are somewhat rare and were not detected in a slightly earlier
study \citep{zeh2015}. Additionally, some authors suggest these data
are impossible to reconcile with some field relationships, challenging
the reliability of U/Pb zircon geochronology for investigating layered
mafic intrusions \citep{latypovchistyakova2022}. Nevertheless, U/Pb
zircon geochronology successfully yields narrow
age ranges for other igneous rocks when these are geologically expected, such as
for erupted ignimbrites \citep{wotzlaw2013}.

\subsection{Origin of Enigmatic Monomineralic Igneous Rocks}
\label{sec:ch5-origin-of-enigmatic-monomineralic-igneous-rocks}

The RLS contains enigmatic monomineralic igneous rocks that continue to
perplex igneous petrologists. For instance, the RLS contains dozens of
centimeter to meter scale, near-pure chromitites which are challenging
to explain. One common explanation involves the mixing of primitive and
evolved magmas inducing chromite saturation \citep{irvine1977}. However, an
in-situ crystallization hypothesis is problematic because the Lower and
Critical Zones contain significantly more chromium than could have been
delivered by a plausible amount of chromium-saturated silicate melt
\citep{eales2000}, and the cumulate rocks above and below some chromitites
are petrologically identical \citep{mondal2007}, seeming to
preclude in-situ magma mixing unless over 300,000~km\textsuperscript{3} of differentiated
mafic lavas were lost from the chamber\textquotesingle s top \citep{cawthornwalraven1998}. The most popular alternative explanation posits
that chromite crystallized within staging chambers and was injected into
the RLS as chromite-phyric silicate magmas \citep{mondal2007,eales2000}. Experiments have shown that chromitites can form through
magmatic sedimentation from a chromite-bearing magma within realistic
timescales \citep{manoochehri2014}.

Stratiform anorthosites within the RLS pose another challenge to
conventional explanations because plausible RLS magma compositions are
rarely, if ever, saturated in plagioclase alone \citep{latypov2020,maier2020}. Various hypotheses propose that these anorthosites
resulted from {[}1{]} density-driven mechanical sorting of a magma
concomitantly saturated in multiple phases \citep{cawthorn2009}; {[}2{]} depressurization-forced saturation of plagioclase alone,
followed by crystal settling \citep{yuan2017} or in-situ growth on
the magma chamber floor \citep{latypov2020}; {[}3{]}
tectonically-induced mechanical sorting of a preexisting cumulate pile
\citep{maier2016}; {[}4{]} reactive porous flow of acidic fluids
\citep{maier2020,sonnenthal1992}; or {[}5{]} restites from partial
melting at the contact of new intrusions of magma \citep{mitchell2019}.

\subsection{The Origin of the Merensky Reef and PGE Mineralization}
\label{sec:ch5-the-origin-of-the-merensky-reef-and-pge-mineralization}

The origin of the Merensky Reef and PGE mineralization in the upper
Critical Zone remains contentious. One theory suggests that immiscible
sulfide fluids scavenged PGEs from a large magma chamber \citep{campbell1983}. Alternatively, percolating fluids rich in chlorine, carbon,
and sulfur may have redistributed PGE elements within an
already-solidified Critical Zone \citep{boudreau1986,ballhaus1985}, or magmatic differentiation controlled by chromite may
have concentrated PGEs within this zone \citep{scoon1994}. The
platiniferous Merensky Reef, a two-part pyroxenite and pegmatoidal
horizon, is debated to be either a composite of multiple synmagmatic or
out-of-sequence sills \citep{mitchell2007,mitchell2019}
or a mixture of disequilibrated crystals precipitating synchronously
from overlying density-stratified Main and Critical Zone magmas
\citep{seabrook2005}. Most authors agree that the Merensky Reef
represents a mixture of Main Zone and Critical Zone magmas \citep{kruger1994,sharpe1985,cawthorn2015}.

\subsection{The Transformation Hypothesis}
\label{sec:ch5-the-transformation-hypothesis}

\citet{vanbiljon1949} developed a controversial hypothesis for the origin of
the RLS: that it represented an extreme case of in-situ metasomatism of
the upper Pretoria Group sedimentary rocks of the Transvaal Supergroup
on the basis of field relationships and geochemical arguments. Indeed,
terminology and concepts from sedimentology and stratigraphy are often
(though arguably erroneously) employed to describe features within the
RLS and other LMIs. Van Biljon's Transformation Hypothesis sought to
explain the origin of the enormously laterally continuous markers within
the RLS, the apparent problem of an enormous amount of displaced
sedimentary rocks, and the transgressive nature of the RLS's contact
with the Pretoria Group. \citet{vanbiljon1949} argues that chert is
transformed to anorthosites, carbonates are transformed to
magnetitites/chromitites, variably impure quartzites are transformed to
(gabbro)norites, and shales are transformed to pyroxenites. He presents
geochemical evidence that seems to show a chemical continuum in
different oxide spaces between unmetamorphosed sedimentary rocks,
variably hornfelsized sedimentary rocks, and mafic rocks of the RLS.

\citet{button1976} emphasized Van Biljon's recognition of the ``space
problem'' introduced by the transgressive RLS into relatively undeformed
sedimentary rocks, instead arguing that the felsic Bushveld rocks
represented the displaced (and transformed) sediments. However, \citet{cawthorn1998}
showed through a simple geometric analysis that the apparently
missing sedimentary rocks may have been preserved underneath the RLS
given the nature of its transgressive and discordant contact with the
Pretoria Group.

One simple counterpoint to Van Biljon's transformation hypothesis for
the RLS is that similar LMIs exist that are emplaced entirely or nearly
entirely within non-sedimentary granitic, basaltic, and gneissic country
rock, such as the chromitite-bearing Great Dyke of Zimbabwe \citep{worst1958} and the finely-layered Skaergaard Intrusion \citep{mcbirney1996}.
Furthermore, the RLS itself onlaps granite in the North Limb \citep{vandermerwe1976}.

Generally, available stable and radiogenic isotopic data require
whole-intrusion-scale magmatic and hydrothermal processes to be the
dominant processes forming the RLS, and for local contamination with
crustal rocks to play a subordinate role \citep{zirakparvar2014,schiffries1989,schiffries1990,willmore2002,
pennistondorland2012,kruger1994,richardson2008}. This general statement has notable exceptions: 
\begin{enumerate} 
    \item The Marginal Zone contains abundant metapelitic migmatites and calc-silicates (\citealp{johnson2003,vanbiljon1949,button1976,frick1967}, 
        and Chapter~\ref{ch:draaikraal} of this work).
    \item Interstitial phases such as anhydrite, almost certainly derived from sedimentary rocks, are documented within the ultramafic
        Lower Zone \citep{wilson2017,yudovskaya2021}.
\end{enumerate}

These observations together suggest that in-situ igneous sediment
transformation is not an essential or dominant process in the formation
of the RLS or other LMIs, but that it may still occur at small scales.
Evidence for this process would become obscured and eventually destroyed
by episodic melting and/or reaction with an abundance of incoming magmas
and hydrothermal fluids, suggesting that this process could also have
been locally important at higher levels within the RLS even when
evidence is subtle or completely lost to igneous overprinting.
